\section{Y-OS and yunix}
\label{sec:os}

\textsf{Y} has a second compilation target beside the hosted Linux ELF:
\ic{ycc --target=kernel} compiles a program that boots on bare metal, and \emph{Y-OS}
is a small operating system built from such programs. It is a target, not a
mode on top of Linux: \ic{res\_target} becomes \ic{"kernel"} and everything
downstream that would otherwise assume an operating system changes behind one
predicate, \ic{target\_is\_kernel()}. A hosted compile is untouched
(Section~\ref{sec:os-cost}).

Y-OS is written in \textsf{Y} alone. The kernel proper is the standard units
\ic{kernel}, \ic{kmem}, \ic{yfs} and \ic{fd}, which travel in the std pack
like every other unit, so a lone \ic{ycc --target=kernel} in an empty directory
builds a bootable kernel with no checkout; the \emph{system} on top of it is
\ic{os/kernel/*.y}, the shell and its programs in \ic{os/uland/}, and the
host-side image builder \ic{os/tools/mkfs.y}.

\subsection{The kernel target}

\begin{lstlisting}
$ ./ycc --target=kernel prog.y -o prog.elf
$ qemu-system-x86_64 -kernel prog.elf -serial stdio -display none
\end{lstlisting}

\noindent The output is an ELF32 Multiboot-1 payload: the entry is 32-bit
code that builds paging, sets PAE/LME/PG and far-jumps into long mode,
because that is the mode a Multiboot loader starts an image in --- the one
step the in-tree assembler cannot express, whose encoder is x86-64 only. The
image is linked at 1\,MB, the serial port COM1 is its output, and \ic{print}
and \ic{puts} work there unchanged (\ic{y\_os\_write} is replaced by a COM1
writer). There is no libc, no Linux and no \ic{/bin}.

\paragraph{What a kernel can use.} The runtime is the same prelude with the
few blocks that reach an operating system replaced: the allocator bumps
through a kernel heap, \ic{y\_os\_read} blocks on the 16550, \ic{y\_os\_exit}
halts the machine, and the five trap stubs print and stop. The string, array
and float cores are pure computation and are used as they are, and the tree
shaker drops what the program does not call exactly as it does for a hosted
program.

A kernel refuses the std units that are one Linux system call away from
working --- \ic{fs}, \ic{net}, \ic{exec}, \ic{prelude}, \ic{crypto}, \ic{x509},
\ic{tls}, \ic{http}, \ic{tui}, \ic{args}, \ic{random}, \ic{test}, \ic{img},
\ic{audio}, \ic{ipc} and \ic{snd} --- each with a one-line reason
(\ic{kernel\_std\_why}), because ``not available'' without one reads as a
compiler limitation. The rule is a deny list, not an allow list: a unit that
is plain computation (\ic{json}, \ic{regex}, \ic{math}, \ic{struct}, \ic{utf8})
is available in a kernel the day it lands. The reverse rule is \ic{std/kernel},
which is \ic{ndef}s into assembly only a kernel image has, so a hosted compile
refuses \ic{import kernel} and says so. Both rules fire only when the unit
resolved \emph{is} the standard one: a program's own \ic{kernel.y} or \ic{fs.y}
beside it shadows std and stays its own business.

\paragraph{\ic{irq N def}.} A top-level declaration that exists only in kernel
mode, whose \ty{N} is the 8259 line (0--15), not the IDT vector:

\begin{lstlisting}
import kernel;

irq 4 def serial_rx() -> void {      // 4 is the 8259 line
    let c = kernel.getc();
    kernel.putc(c);
};
\end{lstlisting}

\noindent The parser records the declaration; the checker requires no
parameters and a \ic{void} result (the generated stub is entered from the
interrupt with nothing to pass and nowhere to put an answer) and refuses a
handler nested inside another \ic{def}; the emitter makes the handler a
body-shake root, because the stub that calls it is generated text the shaker
cannot see; and codegen writes, per declared line, a stub that saves the
caller-saved registers, aligns the stack, calls the handler, acknowledges the
8259 and \ic{iretq}s, and unmasks the declared line. Two handlers on one line,
or one name on two lines, is an error.

The CPU's exceptions (vectors 0--31) get stubs of their own rather than the
halting default: each reports the vector, the error code and, for a page
fault, CR2. A fault is not resumed from.

\subsection{The image and the filesystem}

Y-OS is built as one kernel image and one filesystem image:

\begin{lstlisting}
$ make os            # os/build/os.img and os/build/os.fs
$ make run-os        # boot it here, the serial console on this terminal
$ make run-os-gui    # the same with a screen, keyboard and mouse
\end{lstlisting}

\noindent \ic{make os} needs \ic{qemu-system-x86\_64}. The filesystem is a
YFS1 image; \ic{os/tools/mkfs.y} builds it, \ic{os/manifest.txt} lists what
goes into it and \ic{os/kernel/boot.y} is the kernel program. qemu hands the
kernel one Multiboot module (the filesystem) and it is mounted where it lies:
the files the shell makes and removes live in RAM and are gone when the
machine stops.

\paragraph{YFS1.} Blocks are a superblock field, so the reader is not
hard-wired to a size. Block 0 is the superblock, then the inode table
(64 bytes per inode), then data. An inode has a mode (0 free, 1 file, 2
directory), a size --- bytes for a file, \emph{entry count} for a directory
--- twelve direct block slots and one single-indirect block. A directory entry
is a length, a type, an inode number and a name padded to 4 bytes. An image
must be zeroed where nothing was written, because a free inode is found by
reading a zero mode.

\ic{boot.y} mounts the filesystem, loads \ic{/bin/init} (or, failing that,
\ic{/sh}) into a fresh address space and starts it; when it halts the machine
or exits, so does the kernel. \ic{init} sets the environment
(\ic{HOME}, \ic{USER}, \ic{PATH}, \ic{TERM}, \ic{XDG\_RUNTIME\_DIR}), makes
\ic{/tmp}, \ic{/run/user/1000} and \ic{/home/user}, starts the audio server
\ic{/bin/pulsed} if there is one, runs \ic{/bin/sh} in \ic{/home/user}, and
stays to reap orphans.

\subsection{yunix: the Linux ABI}

\emph{yunix} is the layer that makes the compiler's runtime work on that
kernel: the programs are ordinary \ic{ycc} programs because the kernel
implements the Linux system-call interface the runtime already speaks.
\ic{os/kernel/*.y} are modules of the kernel program \ic{boot.y}, and the
system call table is \ic{sys.y}: \ic{std/kernel.y}'s dispatcher calls
\ic{y\_os\_sys} first for every number, and an answer of
\ic{kernel.SYS\_FALLBACK} hands the call to the kernel's own built-in handler
--- what worked before yunix --- so a stub keeps today's behaviour until its
owner replaces it.

\begin{center}
\begin{tabular}{@{}>{\raggedright\arraybackslash}p{0.24\linewidth}@{\quad}>{\raggedright\arraybackslash}p{0.68\linewidth}@{}}
\toprule
module & what it owns \\
\midrule
\ic{fdt} & per-process descriptor tables: \ic{read}, \ic{write}, \ic{close},
  \ic{lseek}, \ic{ioctl}, \ic{pread}/\ic{pwrite}, \ic{readv}/\ic{writev},
  \ic{poll}, \ic{select}, \ic{dup}/\ic{dup2}/\ic{dup3}, \ic{fcntl} \\
\ic{pipe} & \ic{pipe}, \ic{pipe2} \\
\ic{proc} & \ic{fork}, \ic{vfork}, \ic{execve}, \ic{wait4}, signals (they
  terminate: there are no handlers), ids, \ic{uname}, \ic{kill}, rlimits,
  \ic{reboot} (halt) \\
\ic{time} & a 1\,kHz clock, \ic{clock\_gettime}, \ic{clock\_nanosleep},
  \ic{gettimeofday}, \ic{nanosleep}, \ic{sched\_yield} \\
\ic{thread} & \ic{clone} and \ic{clone3}, \ic{futex}, \ic{arch\_prctl},
  \ic{set\_tid\_address} \\
\ic{mm} & region-based \ic{mmap}, \ic{mprotect}, \ic{munmap}, \ic{brk},
  \ic{madvise}, \ic{getrandom} (\ic{mremap} answers \ic{ENOSYS}) \\
\ic{vfs} & directories and paths over \ic{yfs}: \ic{open}, \ic{openat},
  \ic{stat}/\ic{fstat}/\ic{lstat}/\ic{statx}, \ic{mkdir}, \ic{rmdir},
  \ic{unlink}, \ic{rename}, \ic{chdir}/\ic{getcwd}, \ic{readlink},
  \ic{getdents64}, \ic{truncate}, time stamps \\
\ic{tty} & termios, canonical and raw modes, \ic{VMIN}/\ic{VTIME}, \ic{ISIG},
  window size and job control, over the console \\
\ic{unix} & AF\_UNIX sockets: \ic{socket}, \ic{bind}, \ic{listen}, \ic{accept},
  \ic{connect}, \ic{send}/\ic{recv}, \ic{socketpair}, options \\
\ic{pci}, \ic{ac97}, \ic{dsp} & an AC97 sound card behind \ic{/dev/dsp} with
  the OSS ioctls \\
\ic{tty}, \ic{kbd}, \ic{input} & the console and the input devices \\
\ic{kstat} & the counters the resource monitor and \ic{osstat} read \\
\bottomrule
\end{tabular}
\end{center}

\noindent The hooks the modules plug into live in \ic{std/kernel.y}: the
syscall entry \ic{y\_os\_sys}, the lifecycle calls \ic{y\_os\_task\_forked},
\ic{y\_os\_task\_exec}, \ic{y\_os\_task\_exited} and \ic{y\_os\_task\_switch},
and \ic{block\_restart}/\ic{wake\_chan}/\ic{wake\_expired}. A system call that
must wait does not park in the kernel: it rewinds the user instruction pointer
over the \ic{syscall} and runs again when its channel is woken, so a blocking
read is a normal function that returns when there is something to return.

\paragraph{Memory.} A user program is a Linux executable linked at
\ic{0x400000}, so that address is free in a process's address space. The
kernel maps one 2\,MB page for itself (image, stack, IDT, all below 2\,MB) and
reaches every other physical byte through a direct map at
\ic{0xFFFF888000000000}. Every address space is the kernel's low half, the
direct map, and the process's own user half. An anonymous \ic{mmap} only
reserves a region: the page-fault stub for vector 14 maps a zeroed page on
first touch and resumes the instruction, so a fork copies the few pages a
program touched rather than the whole heap.

\paragraph{The shell and the programs.} \ic{sh} reads a line itself on a
terminal (the tty is raw for the length of the line): cursor and history
editing, Tab completion, quotes and \ic{\bs}, \ic{\$NAME} and \ic{\$?}
expansion, \ic{|}, \ic{<}, \ic{>}, \ic{>>}, \ic{2>}, \ic{;}, \ic{\&\&},
\ic{||}, \ic{\&} and \ic{\#}. Built in are \ic{cd pwd export unset exit env
help true false todesktop} and \ic{NAME=value}; anything else is looked up in
\ic{\$PATH}. Every job runs in its own process group and owns the console
while it runs, so Ctrl-C stops the job and not the shell. Programs in
\ic{/bin} include \ic{ls}, \ic{cat}, \ic{echo}, \ic{mkdir}, \ic{rmdir},
\ic{rm}, \ic{cp}, \ic{mv}, \ic{touch}, \ic{wc}, \ic{head}, \ic{tail},
\ic{pwd}, \ic{stat}, \ic{env}, \ic{clear}, \ic{sleep}, \ic{uname},
\ic{halt}, \ic{write}, \ic{osbench} and \ic{osstat}; a program that reads a
file reads standard input when given none.

\subsection{The screen}

With a display device the machine draws a text terminal and takes keys from a
keyboard and a mouse (\ic{make run-os-gui}, \ic{OS\_VGA=virtio} for
virtio-gpu and virtio-input, or \ic{OS\_VGA=std} for standard VGA with a PS/2
keyboard and mouse). The serial console keeps working beside it, and the
backends follow the host window: a resize gives the terminal new rows and
columns.

\ic{vt.y} is an xterm-style terminal in the kernel (UTF-8 in an 8x16 font,
the 16 colours, a 256-colour cube and 24-bit colour, cursor movement, scroll
regions, the alternate screen, DEC line drawing, mouse reports). The desktop
is what the machine shows at boot: \ic{F1} or \ic{F12} switches between the
terminal and the GUI and the terminal keeps its state while the GUI is up.
\ic{/bin/desktop} shows a wallpaper and a dock whose icons are one per program in
\ic{/bin}; a click types its name into the shell and switches to the terminal.
A file explorer and a resource monitor are panels of the desktop itself. The
wallpaper is a vector image (\ic{os/assets/wallpaper.yvg}) filled by a small
renderer once per size, never per frame. The display is reached through
\ic{/dev/fb0}, which is ioctl-only: information, blit, switch view, type text
and a statistics snapshot.

\subsection{Building, running and testing}
\label{sec:os-cost}

The kernel and system suites are in \ic{make test}:

\begin{itemize}[nosep]
\item the \ic{kernel} suite (\ic{test/kernel/}) builds each fixture with both
  compilers, compares the images, boots them under qemu and checks the serial
  output and the exit status;
\item the \ic{os} suite builds Y-OS with both compilers and requires the
  images to be identical, then boots scripted sessions
  (\ic{test/os/*.cmds}) under qemu and compares the whole serial transcript
  with \ic{test/os/*.expected}. The display has sessions of its own
  (\ic{test/os/display/*.steps}), where each screendump is compared by
  SHA-256 and a mismatch leaves the pictures in \ic{\_check/os-display/}.
\end{itemize}

\noindent \ic{ytest os-session CMD...} scripts a session by hand, and
\ic{YBLESS=1 ./ytest os} rewrites the display hashes after a change that is
meant to move pixels.

\subsection{Limits}

Y-OS is a proof of concept and says where it stops. The kernel heap is a bump
allocator that never gives memory back (64\,MB on a machine with 256\,MB or
more, about 600k path system calls; below that, 8\,MB). Only the 8259
interrupt path is wired: no APIC, no IOAPIC, no MSI. There are no signal
handlers (a signal terminates the task) and no \ic{mremap}. The filesystem and
the terminal live \emph{in} the kernel rather than in user-space servers. The
physical memory allocator hands pages back, but the kernel's own bookkeeping
does not. The GUI has no windows, the screen has no text selection or clipboard
and no font beyond the 8x16 one, and there is no disk: the filesystem is the
boot module. The shell has no job control beyond one foreground job per line
--- no \ic{fg}/\ic{bg}/\ic{jobs}, no Ctrl-Z --- and no shell functions,
\ic{if}/\ic{for}, globbing, history search or wrapped-line editing.
\ic{ycc} itself does not run on Y-OS yet.
Finally, \ic{-g} is refused for a kernel image, because the debug writers are
ELF64/DWARF64.

\paragraph{Cost to a normal compile.} Nothing observable: it is all
behind \ic{target\_is\_kernel()}, the two engines stay byte-identical on the
whole hosted corpus, and the added encoder arms sit at the end of the
instruction ladder, so a hosted line pays for them only if it reaches the end
without matching.
